% titlepage.tex -- титульный лист технического задания.
% Свёрстан по образцу vkr/tz/00-title.tex (шапка капсом, шифр документа,
% объём в листах, блоки подписей на tabular*), с одним сознательным
% отступлением: рамка по периметру листа не рисуется. По п. 2.9 методического
% пособия СЛИ для направления 09.03.02 текстовые документы оформляются на
% листах без рамок; в vkr рамка осталась от бланка другой кафедры.
%
% Лист работает в двух сборках и должен быть верен в обеих:
%   * main.tex          -- ТЗ отдельным документом;
%   * main-appendix.tex -- ТЗ приложением Б к пояснительной записке. Тогда
%     определён \tzappendixletter, и обозначение приложения печатается на
%     этом же листе -- отдельная страница-разделитель под него не заводится.

% ---------- Поля титульного листа ----------
% Титульный лист набран по бланку и правому полю 15 мм, принятому для
% основного текста, не подчиняется: графы подписей заданы в сантиметрах по
% бланку и рассчитаны на полосу 170 мм. Поэтому на время титульного листа
% правое поле остаётся прежним (10 мм), а сразу после него геометрия
% возвращается к основной.
\newgeometry{left=30mm,right=10mm,top=20mm,bottom=20mm}

% Линия для подписи -- линейка, а не \underline: п. 2.13 пособия запрещает
% подчёркивание. Пояснение под линией той же ширины, чтобы стояло по центру.
\newcommand{\tzsignline}{\rule{4.2cm}{0.5pt}}
\newcommand{\tzsigncap}{\makebox[4.2cm][c]{\footnotesize (подпись, дата)}}

% ---------- Объём документа в листах ----------
% Печать односторонняя, поэтому листов ровно столько же, сколько страниц:
% последняя страница минус титульная плюс один. Титульная -- это
% \firstpagenumber минус один (main.tex выставляет счётчик сразу после этого
% листа), последняя берётся меткой LastPage в конце main.tex.
%
% Метку последней страницы даёт пакет lastpage: он ставит её после финального
% \clearpage, когда выведены и отложенные плавающие рисунки приложения (своя
% метка перед \end{document} насчитала бы на четыре листа меньше). Первую
% страницу меткой не отметить -- окружение titlepage сбрасывает счётчик в
% единицу, и она ушла бы в .aux с номером 1, поэтому берём \firstpagenumber.
%
% \getpagerefnumber (пакет refcount) раскрывается, поэтому годится внутри
% \numexpr, в отличие от обычного \pageref. На ещё не известной метке он
% возвращает ноль -- так бывает на первом проходе, latexmk пересоберёт. Но
% так же выглядела бы и потерянная метка, поэтому ноль печатается как «??»:
% на первом же просмотре PDF это видно глазом.
\newcommand{\tzsheets}{%
  \ifnum\getpagerefnumber{LastPage}>0 %
    \the\numexpr\getpagerefnumber{LastPage}-\firstpagenumber+2\relax
  \else
    ??%
  \fi}

\begin{titlepage}
\begin{spacing}{1.0}
\centering
\vspace*{-\baselineskip}

% Обозначение приложения -- наверху посередине листа (п. 12.3 пособия), на том
% же листе, что и титул документа-приложения; следующей строкой в скобках его
% характер (п. 12.5). Техническое задание -- приложение обязательное: оно
% определяет предмет работы, а не поясняет её справочно.
% В самостоятельной сборке макрос не определён, и блок не печатается.
\ifdefined\tzappendixletter
    {\bfseries ПРИЛОЖЕНИЕ \tzappendixletter\par}
    (обязательное)\par
    \vspace{1.5\baselineskip}
\fi

МИНИСТЕРСТВО НАУКИ И ВЫСШЕГО ОБРАЗОВАНИЯ\\
РОССИЙСКОЙ ФЕДЕРАЦИИ

\vspace{\baselineskip}
СЫКТЫВКАРСКИЙ ЛЕСНОЙ ИНСТИТУТ (ФИЛИАЛ) ФЕДЕРАЛЬНОГО\\
ГОСУДАРСТВЕННОГО БЮДЖЕТНОГО ОБРАЗОВАТЕЛЬНОГО\\
УЧРЕЖДЕНИЯ ВЫСШЕГО ОБРАЗОВАНИЯ\\
«САНКТ-ПЕТЕРБУРГСКИЙ ГОСУДАРСТВЕННЫЙ ЛЕСОТЕХНИЧЕСКИЙ\\
УНИВЕРСИТЕТ ИМЕНИ С.\,М.\,КИРОВА» (СЛИ)

\vspace{\baselineskip}
ТРАНСПОРТНО-ТЕХНОЛОГИЧЕСКИЙ ФАКУЛЬТЕТ

\vspace{\baselineskip}
КАФЕДРА ИНФОРМАЦИОННЫХ СИСТЕМ И ТЕХНОЛОГИЙ

\vspace{2\baselineskip}
{\fontsize{14}{16}\selectfont\bfseries ТЕХНИЧЕСКОЕ ЗАДАНИЕ\par}

\vspace{\baselineskip}
к технологической (проектно-технологической) практике.\\
Web-технологии

\vspace{2\baselineskip}
\begin{small}
Тема: разработка CMS для парикмахерских и барбершопов «Сайтама»

\vspace{0.5\baselineskip}
ПП.ТТФ.О.09.03.02.3.230176.ПЗ

\vspace{0.5\baselineskip}
на \tzsheets\ листах
\end{small}

\vspace{2\baselineskip}
\raggedright
\begin{tabular}{@{}p{7.6cm}@{\hspace{0.3cm}}p{4.2cm}@{\hspace{0.3cm}}p{4.4cm}@{}}
    Выполнил: студент 3 курса     &              &                   \\
    очной формы обучения          &              &                   \\
    направления подготовки        & \tzsignline  & Решетняк В.\,А.   \\
    бакалавриата «Информационные  & \tzsigncap   &                   \\
    системы и технологии»         &              &                   \\
\end{tabular}

\vspace{\baselineskip}
\begin{tabular}{@{}p{7.6cm}@{\hspace{0.3cm}}p{4.2cm}@{\hspace{0.3cm}}p{4.4cm}@{}}
    Руководитель:                 & \tzsignline  & Оганезова Н.\,А.  \\
    к.\,э.\,н., доцент            & \tzsigncap   &                   \\
\end{tabular}

\vfill
\centering
Сыктывкар 2026
\end{spacing}
\end{titlepage}
\restoregeometry
